%==============================================================================
%  TRB 2027 投稿草稿 —— 连续动作空间下投影式安全盾的可证明 COLREGs 方向合规无人船避碰
%
%  ▶ 版本：0626-paper.tex（首版中文草稿·基于 paper.tex 原版 TRB/TAD 模板）
%  ▶ 规矩（user 2026-06-26）：
%     · paper.tex = 原版模板·永不改（备份/格式基准）。
%     · 本文件 = 当前工作草稿。先写【中文】，待内容定稿再译英文投稿。
%     · 重大转折/大改 → 复制本文件、按日期命名 + 文件名注明改了什么（如
%       0701-paper-换provable档位B.tex），旧稿即备份，形成时间线。
%  ▶ 编译：因含中文，须用 xelatex（非 pdflatex）：xelatex 0626-paper.tex（跑两次更新交叉引用）。
%
%  ▶ 诚实标注（2026-06-26 撰稿时状态·`02`/`03` L89）：
%     · 主贡献=投影盾把【可证明 COLREGs 方向合规】带进【连续动作空间】（档位A·经验性）。
%     · "provably 零碰撞 / 前向不变性"=档位B/Phase 4·未做 → 标题/正文只敢说"方向合规"。
%     · 结果章为【骨架+占位】：训练在跑、Φ_xtrack 有效性未证、修法A n=5 显著性未达——
%       绝不填未定数字；定稿前所有 [待数据] 须换成多种子 IQM+自助法CI+失败率（不挑种子）。
%==============================================================================
\documentclass[12pt,letterpaper]{article}

%--- 字体：拉丁/公式用 Times 风格；中文用 xeCJK（须 xelatex 编译）---------------
\usepackage{iftex}
\usepackage{amsmath}          % 数学：\text、\arg\min 等（须在字体包前·标准·不改版式）
% 🆕 2026-07-29：0626 模板只载了 amsmath ⟹ \varnothing / \square / \bm 全是未定义命令。
%   本稿的四条命题要用到（空集 / 证毕方块 / 粗体向量）⟹ 补齐。须在字体包前载。
\usepackage{amssymb}          % \varnothing（空集·兜底集为空）· \square（证毕）
\usepackage{bm}               % \bm（粗体向量·速度矢量 \bm{v}）
% 🆕 定理环境（user 2026-07-29 指出"观感不像研究论文"）：原来命题是 \subsubsection，
%   渲染成斜体正文、读者分不清哪里是命题陈述哪里是讨论。改用标准定理环境。
\usepackage{amsthm}
\newtheorem{proposition}{命题}
\newtheorem{lemma}{引理}
\theoremstyle{definition}
\newtheorem*{assumption}{假设}
\theoremstyle{remark}
\newtheorem*{scopenote}{作用域与限制}
\newtheorem*{proofsketch}{证明}
\renewcommand{\proofname}{证明}
\ifPDFTeX
  \usepackage[T1]{fontenc}
  \usepackage[utf8]{inputenc}
  \usepackage{mathptmx}       % 拉丁与公式 Times 风格（pdflatex 分支）
\else                         % XeLaTeX / LuaLaTeX：原生 UTF-8 + 中日韩字体
  \usepackage{mathptmx}       % 先载（提供 Times 数学）；其文本字体族随后被 fontspec 覆盖
  \usepackage{fontspec}
  \usepackage{xeCJK}
  % 🔴 字体按机器自适应（2026-07-29）：原来写死 `Times New Roman` + `Songti SC` 是 **macOS 专有**，
  %    在 Linux（本项目的云端容器、以及多数服务器）上 fontspec 会 **fail-hard、根本编译不出来**。
  %    ⟹ 改成按存在性挑，第一个存在的就用。**macOS 上行为与原来逐字相同**（Times/Songti 排在最前）。
  %    投稿定稿前若要锁死字体，把下面各自的第一候选留下、其余删掉即可。
  \IfFontExistsTF{Times New Roman}{\setmainfont{Times New Roman}}{%
    \IfFontExistsTF{Nimbus Roman}{\setmainfont{Nimbus Roman}}{%
      \IfFontExistsTF{Liberation Serif}{\setmainfont{Liberation Serif}}{%
        \setmainfont{FreeSerif}}}}
  \IfFontExistsTF{Songti SC}{\setCJKmainfont[AutoFakeBold=2.5,AutoFakeSlant=0.2]{Songti SC}}{%
    \IfFontExistsTF{Noto Serif CJK SC}{\setCJKmainfont[AutoFakeBold=2.5,AutoFakeSlant=0.2]{Noto Serif CJK SC}}{%
      \IfFontExistsTF{Noto Sans CJK SC}{\setCJKmainfont[AutoFakeBold=2.5,AutoFakeSlant=0.2]{Noto Sans CJK SC}}{%
        \IfFontExistsTF{Fandol Song}{\setCJKmainfont[AutoFakeBold=2.5,AutoFakeSlant=0.2]{Fandol Song}}{%
          \setCJKmainfont[AutoFakeBold=2.5,AutoFakeSlant=0.2]{WenQuanYi Zen Hei}}}}}
  % 🆕 圈号 ①–⑳ 与全角括号等落在 xeCJK 默认的"西文"类里 ⟹ 会被送给拉丁字体、
  %   而拉丁字体没有这些字形（编译日志报 Missing character）。它们是**九条臂的编号 = 内容**，
  %   ⟹ 显式声明为 CJK 类，走中文字体。
  %   同理中文引号 “ ” ‘ ’（U+2018–U+201F，落在 General Punctuation）也会被送给拉丁字体
  %   ⟹ 渲染成直引号且间距发飘。一并声明为 CJK 类。
  \xeCJKDeclareCharClass{CJK}{"2018 -> "201F, "2460 -> "24FF, "3000 -> "303F, "FF00 -> "FFEF}
\fi

%--- 页面：US Letter，四周 1 英寸页边距，页眉/页脚 0.5 英寸 ----------------------
\usepackage[letterpaper,margin=1in,headheight=15pt,headsep=0.25in,footskip=0.3in]{geometry}

%--- 页眉（斜体 running head）+ 页脚（居中页码）---------------------------------
\usepackage{fancyhdr}
\pagestyle{fancy}
\fancyhf{}
\fancyhead[L]{\itshape 作者姓氏串（Running head·定稿补）}   % ▶ M3：改成本文作者姓氏串
\fancyfoot[C]{\thepage}
\renewcommand{\headrulewidth}{0pt}
\renewcommand{\footrulewidth}{0pt}

%--- 正文：左对齐(右侧自然参差)、单倍行距、首行缩进 0.5 英寸、段间不额外加空 --------
\usepackage{ragged2e}
\setlength{\RaggedRightParindent}{0.5in}
\RaggedRight
\setlength{\parindent}{0.5in}
\setlength{\parskip}{0pt}

%--- 标题层级样式 --------------------------------------------------------------
\usepackage{titlesec}
\setcounter{secnumdepth}{0}                                   % 标题不编号
\titleformat{\section}{\normalfont\bfseries}{}{0pt}{\MakeUppercase}
\titlespacing*{\section}{0pt}{12pt}{4pt}
\titleformat{\subsection}{\normalfont\bfseries}{}{0pt}{}
\titlespacing*{\subsection}{0pt}{12pt}{4pt}
\titleformat{\subsubsection}{\normalfont\itshape}{}{0pt}{}
\titlespacing*{\subsubsection}{0pt}{12pt}{4pt}

%--- 图表标题 ------------------------------------------------------------------
\usepackage{graphicx}
\usepackage{booktabs}
\usepackage{array}
\usepackage{caption}
\captionsetup{labelsep=space,font=bf,justification=raggedright,singlelinecheck=false}
\captionsetup[table]{name=TABLE}
\captionsetup[figure]{name=Figure}

\usepackage[hidelinks]{hyperref}

%--- M1 行号：TRB 硬要求（参考论文每行都带行号）· 0626 稿缺这一项 ---------------
\usepackage[left]{lineno}
\renewcommand{\linenumberfont}{\normalfont\small}

\newcommand{\abshead}[1]{#1\ }
%--- M2 参考文献：TRB 数字顺序引用（正文 (1)、文末 1. …）· 0626 稿用的是 Chicago 作者-年份 ---
\newcounter{refno}
\newcommand{\refitem}[1]{\par\noindent\stepcounter{refno}\hangindent=0.3in\hangafter=1
  \makebox[0.3in][l]{\therefno.}#1\par\vspace{6pt}}
%  正文引用写 \cit{3} 或 \cit{3, 7}；编号须与文末顺序一致（TRB 按首次出现排序）
\newcommand{\cit}[1]{(\textit{#1})}

%--- 占位宏：图未出图前用白色方框占位（图注照最终版写全）----------------------
\newcommand{\figplaceholder}[1]{\fbox{\parbox[c][6cm][c]{\dimexpr\linewidth-2\fboxsep-2\fboxrule}{\centering\textit{【#1 占位 —— 数据到手后替换】}}}}
\newcommand{\TBD}{【待填】}

%==============================================================================
%  TRB 2027 · 中文工作稿  ——  连续动作空间下投影式安全盾的可证明 COLREGs 方向合规无人船避碰
%
%  ▶ 本文件 = 当前工作稿（user 2026-07-29 定名 `729-paper-中文版.tex`）。
%     `_模板_原版勿改_paper.tex` = 原版模板·永不改。
%     `_留档_0626_*.tex` = 四臂时代旧稿·已过时·仅留档，**别照它改**。
%  ▶ 落笔前先读 `写作方案_0729.md`（篇幅预算 / 10 张图表设计 / 数据利用 / 写作顺序 / 结果预案）。
%  ▶ 编译：含中文，须 xelatex（跑两次更新交叉引用）：xelatex 729-paper-中文版.tex
%
%  ▶ 🔴 写作红线（CLAUDE.md §0 · 一个字不能松）：
%     能严证 = 命题1 远场单步无碰 + 命题2 每让路步方向合规 + 命题3 紧急迫近不可避 可靠性证书。
%     绝不写 provably collision-free / 裸 provably compliant / 完整可证明覆盖。
%     命题4 = PROVISIONAL：闭环验证前只写"证了可清障集控制不变 + 设计了终端约束"。
%  ▶ 篇幅：TRB 7500 词上限，图表各折 250 词 ⟹ 10 张图表 = 2500 ⟹ 正文 5000 词。
%     中文稿先往多里写（目标 11000–13000 汉字），定稿转英文时压回。**本会议不允许附录。**
%==============================================================================
\begin{document}
\linenumbers

%------------------------------------------------------------------------------
%  题名与作者信息
%------------------------------------------------------------------------------
\begin{flushleft}
{\bfseries 连续动作空间下基于投影式安全盾的可证明 COLREGs 方向合规无人船避碰}\par
{\itshape（中文工作稿·英文定名待定；标题只敢说“方向合规”，不得出现 collision-free）}\par
\vspace{12pt}

{\bfseries 作者姓名\textsuperscript{*}}\par
所属院系\par
单位名称，城市，国家／地区，邮编\par
ORCID: https://XXX\par
Email: author@university.edu\par
\vspace{12pt}

总页数：\TBD\par                                  % M4：定稿填真数
\vspace{12pt}

\textit{投稿日期：[待定]}\par
\textsuperscript{*}通讯作者\par
\end{flushleft}

\vspace{6pt}

%------------------------------------------------------------------------------
\section{摘要}
\noindent\textit{（结构化摘要·须 < 300 词。“发现”段为占位骨架，数据到手后据真实结果重写，不填未定数字。）}

\justifying

\abshead{研究目标（Objectives）：}可证明 COLREGs（国际海上避碰规则）方向合规的无人船避碰，此前只能二选一：要么把动作空间离散化、屏蔽掉不合规动作，从而放弃连续控制；要么保住连续动作、把 COLREGs 退回成软性奖励惩罚，从而放弃硬保证。本文的目标是把投影式安全盾首次落到海事 COLREGs 的连续动作空间，并精确刻画这一层能证什么、不能证什么。

\abshead{方法（Methods）：}在以安全盾为环境一部分的强化学习框架（SE-RL）下，连续策略（PPO）给出的指令经三步处理：(i) 规则状态机判定相遇态势并给出合规让路方向；(ii) 合规方向半平面、单步无碰条件与动作箱求交，构成安全动作集；(iii) 二次规划把指令投影到该集合上最近的点。我们给出四条命题：远场单步无碰（严格可证）、每让路步方向合规（构造式成立）、紧急迫近不可避的 可靠充分条件证书，以及可清障集的控制不变性（据此设计后备机动终端约束，\emph{结论仍为暂定}）。在 CommonOcean 官方场景库上，以 9 种配置构成的八级单变量递进序列逐级拆开“离散换连续”“无盾换有盾”“两把连续专属改进项”三组机制。

\abshead{发现（Findings）：}\textit{[待数据]} 合规方向作为硬约束进入二次规划 $\Longrightarrow$ 让路步的方向合规不受随机种子方差影响；COLREGs 违规 \TBD；碰撞率 \TBD；油门增量相对离散基线的结构性优势 \TBD；在消融对照中有界动作分布把转艏增量降到 \TBD。所有效率指标报多种子配对检验 + 自助法置信区间 + 达标种子数，绝不挑种子。到达率不列为卖点。

\abshead{创新（Novelty）：}据我们所知，这是首个把投影式安全盾落到海事 COLREGs 连续控制的工作；并给出一套分级的可证明安全层——每条命题都附完整假设与作用域限制，同时明确写出哪些没有被证。

\abshead{实际应用（Practical Applications）：}为无人水面航行器提供规则合规可验证、且保留连续控制平顺性的避碰方案，其保证不依赖训练结果。

\RaggedRight

%==============================================================================
\section{引言}

无人水面航行器正逐步进入与有人船共航的公共水域，其避碰行为必须遵守《国际海上避碰规则》（COLREGs）。与陆上交通规则相比，COLREGs 有一个决定性的不同：它的让路条款是有方向性的。对遇态势中两船各自向右转向；交叉态势中让路船向右避让、直航船保向保速；追越态势中追越船负责让清。也就是说，一个动作即使在几何上避免了碰撞，若转向方向违反规则，仍然是违规的。一条从左侧绕过去、全程没有碰撞的轨迹，在海事上并不合格。

这一点使 COLREGs 合规无法用“最终有没有撞上”来验证，也使它成为安全强化学习中一个不太常见的问题类型：需要保证的不是一个状态量的界，而是每一步动作落在哪个方向的半平面里。现有工作沿两条路推进，各自让掉了一样关键的东西。

\textbf{第一条路是离散动作加动作屏蔽。}\ 把动作空间离散成有限网格，每一步用规则状态机算出合规动作子集，再把不合规的动作屏蔽掉。合规性因此可证——被屏蔽的动作根本不可能被选中，而且这个保证与训练结果无关、与随机种子无关。代价是放弃了连续控制：真实舵机与推进器是连续的，离散网格既引入量化抖动，又把控制精度锁死在网格分辨率上。对本文的任务而言还有一处更具体的损失：COLREGs 要求转向“足够明显”，这在数学上是一个下确界，而网格通常取不到它，只能取到最近的档位，于是每一次合规让路都要多转一些。

\textbf{第二条路是连续动作加软性奖励。}\ 保留连续动作，把 COLREGs 写成奖励函数里的惩罚项。控制是连续的，但保证没有了：奖励只能影响策略的期望行为，不能排除任何一个具体动作。合规率随种子波动，也随训练进程波动。

于是连续动作 $\cap$ 可证明方向合规 $\cap$ COLREGs 这个三重交集是空的。本文用投影式安全盾把它填上：规则状态机给出合规方向的半平面，与单步无碰约束、动作箱一起构成安全动作集，二次规划把策略的连续指令投影到该集合上最近的点。关键在于“可枚举”这个前提可以被“可投影”替代——合规方向恰好是一个半平面，投影到半平面与箱的交集上是一个小规模二次规划，而投影解满足硬约束这件事是构造式的，与训练无关。策略以安全盾在环的方式训练，因而会观察并适应投影带来的修正。

我们同时认为，这类工作真正的难点不只在于把机制搭起来，而在于诚实地说清这一层到底证明了什么。安全强化学习的文献里，“provable”这个词常常被用在一个远小于读者预期的作用域上。因此本文的第二项工作是把可证明层逐条分档：哪一条是严格引理，哪一条是构造式成立，哪一条只是保守的充分条件，哪一条仍然只是设计而未经闭环验证——每一条都附完整假设清单与作用域限制，并明确写出哪些没有被证明。

\noindent\textbf{贡献。}
\begin{enumerate}\itemsep2pt
  \item 连续动作空间中的投影式 COLREGs 安全盾（“\nameref{sec:shield}”一节）。合规方向作为硬区间约束进入二次规划，因此每一个投影可行的让路步的方向合规是构造式成立的，不依赖训练结果、不受种子方差影响。
  \item 一个分级的可证明安全层（“\nameref{sec:provable}”一节）：远场单步无碰（严格可证，且多他船下是无损的合取推广）、每让路步方向合规（构造式）、紧急迫近不可避的 可靠充分条件证书（替换掉此前只检查转向、会漏掉加减速逃逸的不可靠的判据）。我们进一步证明由该证书定义的可清障集是控制不变的，并据此设计后备机动终端约束——\emph{该结论仍为暂定}，闭环接线与官方口径复算尚未完成，本文如实标注。
  \item 一个把机制逐级拆开的实验设计（“\nameref{sec:setup}”一节）：9 种配置构成一条连通的单变量递进序列，相邻两级之间只改动一个变量，从离散基线一路走到完整方法；其中两跳带有不可分割的伴随变量，我们把这一点写在递进序列表的表注里，而不是藏进脚注。
  \item 对边界的完整交代（“\nameref{sec:results}”一节与局限节）：包括少数种子的训练不稳定、追越态势在本基准上无法验证、以及我们不主张的那些东西。
\end{enumerate}

\noindent\textbf{我们不主张的。}\ 不主张可证明无碰；不主张到达率优于离散基线；不主张覆盖机动他船或多船方向冲突；不主张比手工调参的几何控制器更平顺。

%==============================================================================
\section{相关工作}
\textit{[user 2026-07-29 明令：相关工作最后写。参考论文的量级是 33 条参考文献。]}

%==============================================================================
\section{方法}
\label{sec:method}

\subsection{问题设定}

\textbf{本船动力学。}\ 我们忠实沿用对标工作的受偏航约束质点船舶模型（船型 SR108 集装箱船，$175\times25.4$~m）：
\begin{equation}
  \dot x = v\cos\theta,\qquad \dot y = v\sin\theta,\qquad \dot\theta = \omega,\qquad \dot v = a,
  \label{eq:dyn}
\end{equation}
控制约束 $|a|\le a_{\max}=0.24~\mathrm{m/s^2}$、$|\omega|\le\omega_{\max}=0.03~\mathrm{rad/s}$；速度由环境每步裁剪到 $0\le v\le v_{\max}=9.5~\mathrm{m/s}$（步内可短暂达到 $v_{\max}+a_{\max}\Delta t=11.9~\mathrm{m/s}$）；决策步长 $\Delta t=10~\mathrm{s}$，回合上限 170 步（$=1700$~s）。船体矩形的外接圆半径 $R_{\mathrm{circ}}=\tfrac12\sqrt{175^2+25.4^2}=88.4$~m，内切圆半径 $r_{\mathrm{insc}}=\tfrac12\times25.4=12.7$~m。

\textbf{动作空间的口径对齐（关键）。}\ 强化学习策略的常规操作动作箱取 $a\in[-0.048,0.048]$、$\omega\in[-0.018,0.018]$，恰等于离散网格的张成范围（$7\times7=49$ 个网格点，两轴各 7 档）。物理满程 $\pm a_{\max}/\pm\omega_{\max}$ 只留给安全盾与紧急控制器，不属于策略权限。经此处理，连续动作配置与离散动作配置的动作权限完全相同，二者唯一的差别是分辨率（连续区间 vs.\ 7 点网格）——这消除了“连续动作配置因为能开更大舵所以赢”的混淆。离散网格另有一个状态相关的紧急槽位（下标 49），其值由紧急调度器计算，不是网格点。

\textbf{到达判据。}\ 沿用场景库官方的 \texttt{is\_reached}：位置进门 $+$ 朝向在 $\pm9.74^\circ$ 内 $+$ 在时限内。评估侧恒用真门，不放松。

\subsection{投影式 COLREGs 安全盾}
\label{sec:shield}

安全盾在每一决策步依次执行四个环节（图~\ref{fig:overview}a）。

\textit{(i) 态势分类。}\ 规则状态机依据相对方位扇区与碰撞可能性时域，把当前相遇归入
$\rho\in\{\rho_0\ \text{无冲突},\ \rho_1\ \text{直航},\ \rho_2\ \text{对遇},\ \rho_3\ \text{交叉},\ \rho_4\ \text{追越},\ \rho_5\ \text{紧急}\}$，
并在让路态势下给出规则所要求的转向一侧。

\textit{(ii) 构造安全动作集。}\ 规则约束按态势写成动作分量上的区间约束，记为 $U_{\mathrm{colregs}}(\rho)$：
\begin{equation}
  U_{\mathrm{colregs}}(\rho)=
  \begin{cases}
    U_{\mathrm{box}}, & \rho=\rho_0,\\
    \{|a|\le\varepsilon_a\}\cap\{|\omega|\le\varepsilon_\omega\}, & \rho=\rho_1\ \text{（保向保速）},\\
    \{\omega\le-\omega_{\mathrm{turn}}\}, & \rho\in\{\rho_2,\rho_3\}\ \text{（右转让路）},\\
    \{\omega\le-\omega_{\mathrm{turn}}\}\ \text{或}\ \{\omega\ge+\omega_{\mathrm{turn}}\}, & \rho=\rho_4\ \text{（让清一侧由几何决定）},\\
    U_{\mathrm{box}}\ \text{（约束被旁路，见 (iv)）}, & \rho=\rho_5,
  \end{cases}
  \label{eq:ucolregs}
\end{equation}
其中 $\varepsilon_a=0.016~\mathrm{m/s^2}$、$\varepsilon_\omega=0.004363~\mathrm{rad/s}$ 为直航态的保速与保向容差，而
\begin{equation}
  \omega_{\mathrm{turn}}=\Delta_{\text{large-turn}}/T_M=20^\circ/40~\mathrm{s}\approx 0.008727~\mathrm{rad/s}
  \label{eq:omegaturn}
\end{equation}
是“转向幅度足以被对方察觉”这一要求的精确下确界。安全动作集为三者之交
$P=U_{\mathrm{box}}\cap U_{\mathrm{colregs}}(\rho)\cap U_{\mathrm{cf}}$，其中 $U_{\mathrm{cf}}$ 为单步无碰约束，由分离超平面加一阶展开写成线性不等式。

\textit{(iii) 投影。}\ 求解二次规划
\begin{equation}
  u_{\mathrm{safe}}=\arg\min_{u\in P}\lVert u-u_{\mathrm{policy}}\rVert^2
  \label{eq:qp}
\end{equation}
（图~\ref{fig:overview}b）。由于 $P$ 是二维空间中若干半平面与矩形的交，该问题规模极小。

\textit{(iv) 不可行时的兜底。}\ 当 $P=\varnothing$（主要出现在紧急态 $\rho_5$）时，紧急判据在结构上旁路方向约束，兜底控制器转为最小化碰撞风险。这一支是经验性的，不提供保证；“\nameref{sec:prop3}”一节给出它的可判定边界。

策略在 SE-RL 框架下训练，即针对\emph{经过安全盾之后}的转移学习，因而能够观察并适应投影带来的修正，而非训练完成后再由安全盾事后纠正。

\subsubsection{连续控制在合规性上的精细度优势}
式~\eqref{eq:omegaturn} 给出的 $\omega_{\mathrm{turn}}$ 是一个下确界，连续动作可以精确取到它。离散网格的转艏轴取值为 $\{0,\pm0.006,\pm0.012,\pm0.018\}~\mathrm{rad/s}$，其中 $0.006\times40~\mathrm{s}=13.75^\circ<20^\circ$ 并不满足规则要求，故网格上最小的合规档位是 $0.012~\mathrm{rad/s}$，比下确界高出 $37.5\%$。换言之，离散动作屏蔽每执行一次合规让路，转艏幅度都必须超出规则要求约三成七。这一差距是动作空间分辨率带来的结构性差异，与训练过程无关。


%------------------------------------------------------------------------------
\begin{figure}[t]
  \centering
  \figplaceholder{图 1（两格：a 系统框图 · b 动作平面几何）}
  \caption{方法总览。（a）系统框图：观测经连续策略给出指令 $u_{\mathrm{desired}}$，规则状态机判定相遇态势 $\rho$ 与合规让路方向，安全动作集由动作箱、COLREGs 合规半平面与单步无碰约束求交而成，二次规划把指令投影为 $u_{\mathrm{safe}}$ 后送入环境；集合为空时落入兜底链。实线深色框为可证明的模块，虚线浅色框为经验模块——本文对二者的边界不作含糊。（b）动作平面 $(a,\omega)$ 上的几何：灰色矩形为强化学习动作箱（$\pm0.048$, $\pm0.018$），斜线为 COLREGs 合规半平面的边界 $\omega=-\omega_{\mathrm{turn}}$，另一条直线为单步无碰的线性约束，阴影区为可行集 $P$；空心点为离散 $7\times7$ 网格。策略指令落在 $P$ 外时被投影到最近的可行点，箭头长度即投影修正量。可以看到连续动作能精确取到 $\omega_{\mathrm{turn}}$ 这一合规下确界，而网格只能取到最近的档位。}
  \label{fig:overview}
\end{figure}
%  ▶ 来源：手绘矢量图，不由脚本生成；须遵 nature-figure 规范（形制/配色/字号/导出格式）。
%    配色承担语义：可证明块 vs 经验块必须一眼可分，图例直接写"可证明 / 经验"。
%------------------------------------------------------------------------------

\subsection{可证明安全层}
\label{sec:provable}

本节的口径：“证明”一律指严格引理；存在性级、经验级结论逐条标注。完整证明移交外部技术报告（正文无附录）。

\subsubsection{远场态势：单步无碰}

\begin{proposition}[远场单步无碰]
\label{prop:farfield}
设本船与单个他船的中心距为 $d=\lVert p_e-p_m\rVert$。若
\begin{equation}
  d \;\ge\; d_{\mathrm{safe}} + v_{\mathrm{obs,max}}\Delta t + \rho_{\mathrm{ego}},
  \label{eq:farfield}
\end{equation}
则一个决策步之后，两船（以外接圆保守占据）必不相交。其中 $d_{\mathrm{safe}}=R_{\mathrm{circ,ego}}+R_{\mathrm{circ,obs}}$，$\rho_{\mathrm{ego}}=v_{\max}\Delta t+\tfrac12 a_{\max}\Delta t^2$ 为本船单步最大位移。
\end{proposition}

代入常数：$R_{\mathrm{circ,ego}}=88.4$~m；他船宽度不在状态里，故取方形保守上界并叠加安全裕度 $d_{\mathrm{obs,safety}}=2\ell$，得 $R_{\mathrm{circ,obs}}=\tfrac12\sqrt2\,\ell+2\ell=123.74+350=473.74$~m；于是 $d_{\mathrm{safe}}=562.16$~m，$\rho_{\mathrm{ego}}=107.0$~m，$v_{\mathrm{obs,max}}\Delta t=95$~m，阈值 $=764.16$~m（工程取 $\sim780$~m 留裕度）。

\textit{假设。}\ (A1) 单他船（多他船须对每个 $i$ 分别满足）；(A2) 他船速度 $\le 9.5$~m/s——原文只对本船设速度上限，这一条必须显式假设，我们在全池上实测：2000 个基准场景、342{,}000 个障碍状态，最大值 $7.10$~m/s，零个超过 $9.5$（裕度约 $2.4$~m/s）；若换成航速更高的真实 AIS 数据，阈值须相应抬高。(A3) 本船步末 $v\le v_{\max}$（环境裁剪保证）；(A4) 他船步内近似恒速；(A5) 外接圆保守占据。

\textit{证明。}\ 由三角不等式，$d_{\mathrm{next}}\ge d-\lVert\Delta p_e\rVert-\lVert\Delta p_m\rVert\ge d-\rho_{\mathrm{ego}}-v_{\mathrm{obs,max}}\Delta t\ge d_{\mathrm{safe}}$，两外接圆不交，故两船体不交。$\square$

\textit{作用域。}\ 三角不等式与方向无关，故对遇已是真正的最坏情形，几何上没有缺口——这是整个可证明层最稳的一块。它是单步充分条件，不是不变性，也不是近场保证；目前尚未实现为提前退出的快路径。

\subsubsection{让路态势：逐步的方向合规}

\begin{proposition}[每让路步方向合规]
\label{prop:compliance}
在被状态机判为让路态势（$\rho\in\{\rho_2,\rho_3,\rho_4\}$）、且投影二次规划可行（未进入兜底）的每一个决策步上，实际执行的 $u_{\mathrm{safe}}$ 必落在式~\eqref{eq:ucolregs} 所定义的合规方向半平面内。
\end{proposition}

\textit{证明。}\ 合规半平面以硬区间约束进入二次规划，投影解必然满足它。$\square$（构造式，不需要独立定理。）

\textit{假设与作用域。}\ (B1) 单他船；(B2) 状态机分类正确（继承约 $12.6\%$ 的分类差与并列判定的已知缺口）；(B3) 二次规划可行；(B4) 作用域 $=$ 让路态势的相遇。实测让路触发率约 $13\text{--}16\%$，且这是\textbf{相遇（回合）级}而非\textbf{逐步级}的比率；直航违规才按步计。

\noindent\textbf{作用域限制。}\ 本命题不覆盖紧急态（占冲突步的 $55\text{--}60\%$，其判据在结构上旁路合规约束），也不覆盖放松与最小碰撞风险两支兜底（它们丢掉半平面）。因此只能写\textbf{"每让路步方向合规"}，\textbf{绝不能写"每步合规"}——实测每局违规数非零。

\subsubsection{紧急态势：碰撞不可避的保守判据}
\label{sec:prop3}

紧急控制器是经验性的最后手段。我们不主张它能避碰；相反，我们给出一个\emph{可靠}（sound）的充分条件——所谓可靠，指该判据只会漏报、\emph{不会}误报——用以判定“碰撞在物理上已不可避”，从而把这个黑箱变成一条可判定的边界。它替换掉此前一版只检查转向逃逸、因而会漏掉加减速逃逸的不可靠的判据。

\begin{proposition}[迫近不可避判据]
\label{prop:imminent}
设单个恒速他船。定义本船\emph{可达中心过近似集} $R_{\mathrm{box}}(t)$（艏向系矩形，中心取恒速外推），纵向半轴 $B_{\mathrm{lon}}(t)=\tfrac12 a_{\mathrm{eff}}t^2$，横向半轴 $L_{\mathrm{lat}}(t)$，其中 $a_{\mathrm{eff}}=\sqrt{a_{\max}^2+(v_{\mathrm{bnd}}\omega_{\max})^2}$、$v_{\mathrm{bnd}}=11.9$~m/s，且
\begin{equation}
L_{\mathrm{lat}}(t)=
\begin{cases}
\dfrac{v_{\mathrm{bnd}}}{\omega_{\max}}\bigl(1-\cos\omega_{\max}t\bigr), & \omega_{\max}t\le\pi/2,\\[1.4ex]
\dfrac{v_{\mathrm{bnd}}}{\omega_{\max}}+v_{\mathrm{bnd}}\Bigl(t-\dfrac{\pi}{2\omega_{\max}}\Bigr), & \text{否则.}
\end{cases}
\end{equation}
设 $O(t)$ 为他船船体矩形。若
\begin{equation}
  \exists\,t^\ast\in[0,T]:\quad R_{\mathrm{box}}(t^\ast)\subseteq\bigl(O(t^\ast)\oplus\mathrm{disk}(r_{\mathrm{insc,ego}}=12.7~\mathrm{m})\bigr),
  \label{eq:imminent}
\end{equation}
则碰撞不可避：任何可行控制序列都会在 $t^\ast$ 与该他船相撞。
\end{proposition}

\textit{证明（三步）。}\ (i) $R_{\mathrm{box}}$ 过近似真可达中心——速度矢量变化率 $\lVert d\bm v/dt\rVert=\sqrt{a^2+(v\omega)^2}\le a_{\mathrm{eff}}$，二重积分给出纵向界，横向偏移由 $\int_0^t v_{\mathrm{bnd}}\sin(\omega_{\max}s)\,ds=L_{\mathrm{lat}}(t)$ 界住；(ii) 内切盘欠近似船体；(iii) 若可达中心全部落在 $O(t^\ast)$ 的 $12.7$~m 膨胀内，则本船内切盘（$\subseteq$ 船体）与他船体相交 $=$ 碰撞，且这对所有可达位置成立。$\square$

\textit{验证。}\ 六路对抗审计（20{,}000 随机 $+$ 735 精选场景，含三段 bang-bang 逃逸搜索与精确船体碰撞几何）零真假阳；把包络硬化之后，基准集加 1500 例模糊测试仍 0 假阳、且判据非空（触发于 $t^\ast\le7$~s）。

\noindent\textbf{覆盖范围。}\ 本判据是可靠的（单侧保守的充分条件）。不是早预警（它是末端分类器，不是规划代价）；不完备（不触发 $\ne$ 可避）；不覆盖多障碍（逐对可避不能合成）；不覆盖机动他船（必须恒速）。\emph{绝不写}“可证明零碰撞”——这只是紧急层的分类器，不是全系统保证。

\subsubsection{多他船的推广}

碰撞自由在数学上是对全部他船的合取，这决定了推广方向：充分条件按 $\forall i$ 合成，不可避证书按 $\exists i$ 合成。

\textbf{命题 1 的推广：}若对所有 $i$ 都满足同一个 $\sim780$~m 阈值，则一步后本船体与任何他船体都不相交；单他船的证明对每个 $i$ 独立成立，可靠且不额外保守。新增义务：每个他船都要满足 $v_{\mathrm{obs}}^i\le9.5$。

\textbf{命题 2 的推广：}设让路他船集合 $G$、各自合规半平面 $H_i$。当 $H_G=\bigcap_{i\in G}H_i\ne\varnothing$ 且二次规划可行时，执行动作对所有 $i\in G$ 方向合规。适用边界比单他船更尖锐：若两条他船要求相反的让路方向，则 $H_G=\varnothing$，该步落到兜底，记为一次违规。这不是本方法的缺陷，而是 COLREGs 自身的边界——它是一套成对规则，对多船没有清晰规范。我们可辩护的立场是：“成对规范相容时合规，不相容时避碰优先”。随 $N$ 增大可行集收缩、让路合规率下降；这必须实测，不能假设。

\textbf{命题 3 的推广：}$R_{\mathrm{box}}$ 与他船无关，故只要存在某个 $i$ 与某个 $t^\ast$ 使包含关系成立，碰撞即不可避。不能合成的那一半是“可避性”：所有 $i$ 都不触发，不蕴含全局可避（躲开 A 可能被逼进 B 的未来位置）。

\subsubsection{可清障集上的前向不变性（暂定结论）}
\label{sec:prop4}

% 🔴🔴 书写纪律（`Paper/README.md` + `02_理论推导/命题4_..._0723.md` 的 v2.1 三条更正）：
%   ① 证的是"可清障集 A 控制不变"，不是"盾已前向不变"——盾还没接上终端约束；
%   ② 残余 ≈ 0（不是 46%、也不是 5%）：门2 的 in-A 条件率三态全 100%，
%      crossing 那 2.9% 是 not-in-A（压根不可清障、落兜底），与合规冲突无关；
%   ③ 别写"每让路步 by construction"——门2 测的是"该状态**存在**合规 certified 后备"。

命题 1--2 都是单步保证。单步无碰不自带递归可行性：一个此刻安全的状态，可能没有任何安全的后续。这也是反应式控制障碍函数滤波器在 $\Delta t=10$~s 这样长的零阶保持下所缺的东西。我们的升级思路如下。

\textbf{可清障集 $A$。}\ 固定单个恒速他船。称有限控制序列 $m$ 是状态 $s$ 的已认证直尾逃逸，若：(i) 它可执行——每段边界落在决策步边界上，因而 $m$ 是逐步恒定控制的序列；(ii) 尾段恒速直行（$\omega=0$ 且 $a=0$）；(iii) 一个可靠的 Lipschitz 净空证书在 $[0,H]$ 上给出船体距离 $>0$，且已过最近点并留有裕度。定义 $A:=\{s:\exists\ \text{已认证直尾逃逸}\}$。

\begin{lemma}[恒速直尾蕴含永久净空]
\label{lem:convex}
若尾段恒速直行且他船恒速，则两个朝向固定、以恒定相对速度平移的矩形之间的距离函数 $g(t)$ 是凸的；因此一旦过了最近点开始增长，便对所有更晚的 $t$ 持续增长。
\end{lemma}

\noindent\textbf{可靠性要点。}\ 凸性要求尾段恒\emph{速}，仅 $\omega=0$ 不够——加速的直行尾段路径对 $t$ 非仿射，$g$ 未必凸；故证书同时要求 $a=0$。该证书是充分条件：可靠，但保守，会拒绝一些其实安全的逃逸。

\begin{proposition}[可清障集的控制不变性]
\label{prop:invariant}
对每个 $s\in A$ 及其已认证逃逸 $m=(u_0,u_1,\dots)$，施加 $u_0$ 一个决策步后得到 $s'$，其尾 $m'=(u_1,u_2,\dots)$ 仍是 $s'$ 的已认证直尾逃逸——因为它就是同一条已通过可靠性证书的物理轨迹的后半段。故 $s'\in A$，即 $A$ 是控制不变的。据此，“后继仍在 $A$ 内”这一终端约束（实现为对一条具体后备机动的 $O(1)$ 重认证）在设计上可使带盾闭环在 $A$ 上递归可行。
\end{proposition}

\textbf{合规版本 $A\cap U_{\mathrm{colregs}}$。}\ 盾必须选一个合规的第一步。设合规后备为“向规定一侧转 $t_1$ 秒、然后直行”。情形 1（$t_1>\Delta t$）：尾段仍以合规转向开头、仍已认证，故 $s'\in A\cap U_{\mathrm{colregs}}$——这一情形已证。情形 2（$t_1\le\Delta t$）：尾段第一步是直行，需从 $s'$ 重新找合规逃逸，一般性证明尚缺。在盾的真实让路状态分布上，情形 2 从未出现（$175\text{--}260$~m 的大型船体、转艏率约 $1.7^\circ$/s，任何合规清障转向都必然超过一个决策步）。

\noindent\textbf{暂定与作用域。}\ (0) 终端约束只在让路后继上认证 $A$ 归属；降级为无冲突的后继平凡安全，直航与紧急后继分别落回命题 1--2 与经验兜底。(i) 单恒速他船；多他船更弱（$A=\bigcap_i A_i$，逐目标的可行域可能互相冲突）。(ii) $A$ 由有限机动族加保守证书判定，是真可清障集的内近似。在真实让路分布上，可清障状态里每一个都存在 COLREGs 合规的已认证脱离机动（三种态势的条件率均为 $100\%$）；唯一的缺口是交叉态势中约 $2.9\%$ 的状态\textbf{压根不可清障}（任何方向都清不掉）而落到兜底——\textbf{这是可清障性的边界，不是"合规与安全冲突"}。(iii) 终端约束模块已实现并通过单元测试，\textbf{接进部署中的盾并验证闭环仍待完成}；上述比率来自高精度本地积分器，官方口径复算待补。(iv) 这些让路状态取自表现良好的带盾回放，故比率刻画的是可清障集的\textbf{几何}，不是对抗性迫近冲突下的性能。

%------------------------------------------------------------------------------
\begin{table}[t]
  \caption{可证明安全层：保证、状态与诚实作用域}
  \label{tab:provable}
  \centering\small
  \begin{tabular}{@{}p{0.20\linewidth}p{0.24\linewidth}p{0.46\linewidth}@{}}
    \toprule
    保证 & 状态 & 作用域与限制 \\
    \midrule
    远场单步无碰（命题 1，含多他船）
      & 已证（初等引理；多他船合取不额外保守）
      & $d\ge\sim780$~m；他船恒速；$\forall i$ 合取；每个 $v_{\mathrm{obs}}\le9.5$~m/s；单步，非不变性；尚未实现提前退出 \\
    \addlinespace
    每让路步方向合规（命题 2，含多他船）
      & 构造式成立（作用域受限）
      & 让路相遇（约 $13\text{--}16\%$，\textbf{相遇级}非逐步级）；不含紧急态与兜底；多他船下随 $N$ 下降；方向冲突 $\to$ 兜底（COLREGs 成对性边界）。\textbf{绝不写"每步"} \\
    \addlinespace
    紧急迫近不可避证书（命题 3，含多他船）
      & 可靠充分条件（$>$20{,}000 对抗场景 0 假阳）
      & 仅迫近；他船恒速；$\exists i$ 可合成、可避性不可合成；不完备、非早预警；是分类器不是系统保证 \\
    \addlinespace
    可清障集 $A$ 上的前向不变 / 递归可行（命题 4）
      & \emph{暂定}：$A$ 的控制不变性已证、证书可靠（长时域模糊 0 假阳）；终端约束模块已单测；闭环接线与官方口径复算待补
      & 单恒速他船；$A$ 是内近似；可清障态中合规已认证脱离机动的条件存在率三态均 $100\%$；唯一缺口是交叉态约 $2.9\%$ 状态\textbf{压根不可清障} $\to$ 兜底（非合规冲突）；$A\cap U_{\mathrm{colregs}}$ 的不变性\textbf{仅情形 1 已证} \\
    \bottomrule
  \end{tabular}
\end{table}
%  ▶ 本表内容已定稿，零待填。数据来源：Paper/02_理论推导/ 全部六份。
%------------------------------------------------------------------------------

我们能辩护的头条主张（挂满假设）：连续动作 $\cap$ 可证明方向合规 $\cap$ COLREGs（此前空白的三重交集），加上远场单步可证明安全与紧急迫近不可避的 可靠性证书。\emph{我们不主张}：裸的可证明无碰、完整可证明覆盖、机动他船、多船方向冲突。

\subsubsection{紧急兜底与失败机理的边界}
当投影在紧急态不可行时，判据在结构上旁路方向约束，兜底控制器最小化碰撞风险——这是经验兜底，不是保证。在一个专门逼出紧急支的对抗性纯几何基线上，我们观察到的碰撞全部发生在紧急归口的步上，即“最后手段尝试并失败”，且都落在命题 1 的适用域（远场、非紧急、单步）之外。它们不是对头条主张的反例，但坐实了紧急态不是保证。\textit{[正式实验重跑后此处数字改为 \TBD，来源见 “\nameref{sec:res-shield}”]}

\subsection{连续动作的两项改进}
\label{sec:keys}

把 COLREGs 盾搬进连续动作空间之后，出现了两个连续动作空间特有的问题，我们各配一项改进。“\nameref{sec:arms}”一节的消融（⑦⑧⑨①）就是用来拆开这两项改进各自贡献的。

\textbf{改进一：有界动作分布。}\ 默认的连续策略用无界对角高斯 $+$ 硬裁剪。实测（直读 10 个存档的对数标准差）：$\sigma_a=1.438\times$ 半箱、$\sigma_\omega=2.182\times$ 半箱；更关键的是，真实观测下策略的确定性均值本身就有 $76\text{--}84\%$ 的步落在动作箱之外（中位数 $2.0\text{--}3.2$ 倍半箱）——这意味着\textbf{即使把探索方差压到零，裁剪也照样把动作按死在满舵上}，压方差、退火熵系数这类便宜修法全部无效。轨迹反推显示 $73.5\%$ 的步贴在箱边、满舵左右一步翻面的占 $23.3\%$（贡献了 $53.0\%$ 的 $|\Delta\omega|$）。机理有两条：(a) 常见实现把\textbf{裁剪后}的动作送进环境、却把\textbf{未裁剪}的动作存进缓冲区，于是均值一旦出箱，再往外挪奖励一分不变 $=$ \textbf{梯度平高原}，平顺度惩罚项无处着力；(b) 高斯的熵 $=\sum\log\sigma+\text{常数}$、\textbf{无上界}，而熵系数是常数，于是标准差有一条"白拿熵奖励"的上行通道。

改用 Beta 分布同时堵死这两条：支撑天生就是闭区间，动作不可能越箱 $\Rightarrow$ 裁剪层退化成恒等 $\Rightarrow$ 平顺度惩罚处处有梯度；且 Beta 的熵有上界（均匀分布），熵系数顶不出去。我们强制 $\alpha,\beta\ge1$，这一条不可省：$\alpha$ 或 $\beta<1$ 时 Beta 是 U 形（两端密度发散）$=$ 数学形式的 bang-bang，而且此时众数落在两个端点、区间内的驻点是密度极小值——由于所有报数都用确定性动作，这会静默污染每一个数字而在聚合指标上完全看不出来。\textit{（注意 $\alpha,\beta\ge1$ 并不阻止动作饱和：实测 $\alpha=8.19,\beta=1.002$ 时众数 $=0.9997$，仍贴着箱边。它只保证单峰与端点密度有限；真正治病的是“支撑全在箱内”与“熵有上界”。）}

\textbf{改进二：对称的让路入口。}\ 对标工作的状态机中，“无冲突 $\to$ 让路”这一转移只认持续性判据，而“直航 $\to$ 让路”另有一支即时判据。我们补上对称的一支：持续性判据未命中时，再看“现在是否已经成立”，成立即进入让路态。这项改进改变盾实际施加的动作，因此带盾配置必须重训，且训练与评估必须同档。

%==============================================================================
\section{实验设置}
\label{sec:setup}

\subsection{场景库与划分}
场景库为 CommonOcean 的双船人工场景库，共 2000 个场景，官方划分为训练 1400 / 测试 600。我们把官方训练 1400 再切成 1300 训练 $+$ 100 验证，训练与训练期评估全程不碰测试集。

\textbf{报数分母 $=600$。}\ 以往探索期实验用的是小场景集，其训练场景与官方测试 600 有交集、须剔除。本次全部配置改在官方 1300 上训练，而官方 1300 与测试 600 零交集，因此没有任何一种配置贡献泄漏，分母就是 600。我们已独立复核了官方划分（自行重跑同一划分并求交集：训练 $\cap$ 测试 $=0$、验证 $\cap$ 测试 $=0$、训练 $\cup$ 验证 $=$ 官方训练 1400、训练 $\cap$ 验证 $=0$）。
\textit{口径提示：600 与探索期的分母不是同一个数，两边的结果不放进同一张表；正文若引用探索期结论会标明那是另一套口径。}

\subsection{九种配置与八级单变量递进序列}
\label{sec:arms}

表~\ref{tab:arms} 列出九种配置。它们构成一条连通的单变量递进序列：相邻两级之间只改动一个变量，因此每一个机制的贡献都能单独读出来。

%------------------------------------------------------------------------------
\begin{table}[t]
  \caption{九种配置的配方与单变量递进序列}
  \label{tab:arms}
  \centering\small
  \begin{tabular}{@{}clllcccc@{}}
    \toprule
    \# & 配置 & 作用 & 动作空间 & 盾 & 塑形 & 分布 & 相对上一级只差 \\
    \midrule
    ① & 本文方法      & 主角                    & 连续 & 投影盾   & 全 & Beta  & （由 ⑧ 或 ⑨ 各补一项改进）\\
    ② & 离散-有盾     & 唯一正面对手（对标）    & 离散 49 & 动作屏蔽 & 零 & ---   & 合规怎么施加（相对 ③）\\
    ③ & 基线          & 上界参照 · 递进序列的根     & 离散 & 无       & 零 & ---   & ——\\
    ④ & 软奖励        & 钉死“软奖励换不来合规”  & 离散 & 软奖励   & 零 & ---   & 合规怎么施加（相对 ③）\\
    ⑤ & 连续-无盾     & 项目至今从未训过        & 连续 & 无 & 零 & 高斯 & 动作空间（相对 ③）\\
    ⑥ & 连续-有盾     & 与 ⑤ 逐字同配方         & 连续 & 投影盾   & 零 & 高斯  & 只差盾（相对 ⑤）\\
    ⑦ & 消融·都不改   & 消融基准                & 连续 & 投影盾   & 全 & 高斯  & 只差连续专属塑形（相对 ⑥）\\
    ⑧ & 消融·只 Beta  & 改进一的贡献            & 连续 & 投影盾   & 全 & Beta & 只差有界分布（相对 ⑦）\\
    ⑨ & 消融·只对称   & 改进二的贡献            & 连续 & 投影盾   & 全 & 高斯  & 只差让路入口（相对 ⑦）\\
    \bottomrule
  \end{tabular}

  \vspace{4pt}
  \footnotesize
  \textbf{表注（三条限定必须与本表同处出现）。}
  (1) ③$\to$⑤ 这一级不是纯单变量：环境层与奖励层是干净的（我们做过数值数值比对，差为 0），两个算法的默认超参也已逐字核对、完全相同；但算法类（带屏蔽的 PPO vs.\ 普通 PPO）、策略族（49 类别分布 vs.\ 2 维连续分布）、熵系数的语义（离散熵与微分熵不是同一个量）三者是“换动作空间”不可分割的伴随。故这一级只能读作“动作空间及其必然伴随的贡献有多大”。
  (2) ⑥$\to$⑦ 的“连续专属塑形”是 4 件一起开的捆绑包（进门势 / 横向进带势 / 终端保速势 / 平顺度惩罚），本次没有逐件拆开 $\Rightarrow$ 只能写“这套组合贡献了 $X$”。
  (3) ⑦$\to$⑧ 换分布带一个必然伴随：两种分布的初始探索尺度在加速度轴上差约 2.5 倍。高斯的初始标准差是一个由较窄的转艏轴动作箱派生的标量（$\sigma\approx0.009$），两轴共用；加速度轴箱宽是转艏轴的 2.67 倍，故高斯在加速度轴上相对更窄（$\sigma/$半宽 $=0.19$ vs.\ 转艏轴 $0.50$）；Beta 的初始形状按箱宽相对定义（$\approx0.48\times$ 半宽，两轴相同）。这是换分布的必然伴随，不是接线错误——不存在“既换分布又保持初始探索尺度不变”的写法，因为高斯的尺度是绝对的、Beta 的是相对的。
\end{table}
%------------------------------------------------------------------------------

\textbf{⑤⑥ 这对配对配置的硬限定（必须与“盾的贡献有多大”这个数字同处出现）。}\ 代码守卫规定无盾的连续动作配置只能用极简配方（无界高斯 / 原文让路入口 / 零塑形），因为让无盾配置开连续专属塑形会造成对离散无盾基线的不对称、污染归因。而无界高斯正是“\nameref{sec:keys}”一节测出 $73.5\%$ 的步贴箱边的那个形态。$\Rightarrow$ \textbf{"盾在连续动作空间里值多少"这个倍数，是在连续动作最差的形态下测得的。}同理，③$\to$⑤ 的"离散换连续本身值多少"\textbf{很可能是负的}——那是诚实的下界，我们照写。

\subsection{训练预算、种子与算力}

每条 run 从零端到端训练 \textbf{10{,}158{,}080 步}（20 段 $\times$ 507{,}904），不用热启动；分段存档全程开（每段留一份 $\Rightarrow$ 20 个存档候选）；训练期评估用那 100 个验证场景，曲线即验证曲线。每条 run 用 8 个并行环境。

种子 $=$ 0--7，共 8 颗，分布在 4 台机器上（每台 2 颗、每台 6 路并发）。\emph{按种子分机器，绝不按配置分机器}——同种子配对比较是本文的主统计工具，必须落在同一台机器内；按配置分机器等于把“机器”变成混淆变量。评估另说：一张报数表里每种配置必须同一台机器、同一轮内评完。

报数预算取所有 run 都达到的最小段数，同一张表全部配置同一个值，并报实际步数而非名义值。

\subsection{指标}
主指标对齐对标工作：到达率、碰撞率、每局 COLREGs 违规数、紧急控制步占比、每局时长；违规数进一步拆成让路违规与直航违规。控制质量：转艏增量（分总计 / 让路步 / 其他步）、油门增量、jerk、控制努力、次网格细调率、速度反转次数。安全裕度：最近点中心距 / 净空 / 发生时刻、最近点处航向误差与速度。盾的行为：介入步数、兜底触发率、投影修正量（均值 / p95 / 最大 / 零修正比例）。分型：按会遇类型拆全部指标、失败原因分型。实时性：单步求解耗时，分盾 / 策略 / 整步三层。

\textit{采集边界（决定了图能怎么画）。}\ 训练期每 507{,}904 步记录验证曲线；每个 rollout（16{,}384 步）记录策略与值函数损失、熵、KL、裁剪比例、可解释方差、每局回报，以及动作饱和率、舵反转率、盾介入率、盾改写量等行为量（转艏与加速度两轴分开记，盾改写量按箱半宽归一化后再合成，不做量纲直加）。训练期没有违规曲线——违规计数只在评估侧接线，因此违规只有那 20 个评估点。另外，各配置奖励口径不同，每局回报不可跨配置比较。

\subsection{存档选取与两轮报数}
选取规则先定死、不带任何旋钮：候选 $=$ 全部 20 段；单一指标 $=$ 验证集到达率，取最高；平局取更早的段。\emph{绝不用测试集挑存档}，所有配置一视同仁。

必须两版都报：验证集 $N=100$，到达率的二项标准误约 4 个点；在 20 个候选里取最大值，期望上抬约 7.5 个点——比我们专门立规矩去防的跨轮次抖动（约 1.6 个百分点）还大 4 倍多。故最佳存档一轮 $+$ 末段存档一轮，两版都报并量化交代这个偏倚。

\subsection{统计方法}
\label{sec:stats}
主检验 $=$ 同种子配对符号检验。8 颗种子时双侧最小可达 $p=2/2^{8}=0.0078$，足以支撑“统计上可辨”的结论。区间用按种子重采样的自助法 95\% 区间；\textbf{碰撞}用 Fisher 精确检验（碰撞是稀有事件，正态近似不适用）。全部自算，不调统计库，脚本随论文给出。

\textit{[若最终达标种子数掉到 3--4，主检验改为按测试场景配对（每种配置在每个测试场景上的逐局明细，$N=600$，样本量比逐种子大一个量级），逐种子那一版降级为稳健性检查。换检验的唯一理由是算力受限，会在正文写明，且逐场景配对不能替代逐种子回答“训练可靠性”——那条只能靠种子，$n$ 小就如实写 $n$ 小。]}

“达标种子数”这一列附判据定义：达标 $=$ 严格到达率 $\ge50\%$；\textbf{崩}（打转吸引子）$=$ 到达 $<50\%$ \textbf{且}末段每局时长 $\ge1360$~s（$=0.8\times$ 回合上限）；\textbf{欠训} $=$ 到达 $<50\%$ 且末段每局时长 $<1360$~s。\textbf{并报过渡带（1100--1600~s）的条数}——那一带的标签本身是灰的。（全库 100 条低到达 run 的实证：崩组落在 $[1377,1700]$、欠训组落在 $[567,1354]$，阈值处不重叠，$12\%$ 落在过渡带。）

%==============================================================================
\section{结果}
\label{sec:results}
\textit{[本节全部为占位。数据到手前\emph{绝不先编任何数字}。每个占位的来源脚本写在紧邻的注释里。]}

\subsection{主表：九种配置的同口径对照}
\label{sec:res-main}

%  ▶ 来源：Paper/正式实验/03_表/make_头条表.py → 表1_头条表_*.md
%    区间与 p 值：Paper/正式实验/05_统计/stat_检验与区间.py
%    前置：两趟重评产出 02_重评产物/g*.json
\begin{table}[t]
  \caption{主表：九种配置在官方测试集上的同口径对照（最佳存档轮 · 全部种子）}
  \label{tab:main}
  \centering\small
  \begin{tabular}{@{}lcccccccc@{}}
    \toprule
    配置 & $n$ & 达标 & \textbf{到达率 \%[CI]} & \textbf{碰撞 \%(局)} & \textbf{违规/局} & \textbf{转艏$\Delta$} & \textbf{油门$\Delta$} & \textbf{紧急步\%} \\
    \midrule
    ① 本文方法      & \TBD & \TBD & \TBD & \TBD & \TBD & \TBD & \TBD & \TBD \\
    ② 离散-有盾     & \TBD & \TBD & \TBD & \TBD & \TBD & \TBD & \TBD & \TBD \\
    ③ 基线          & \TBD & \TBD & \TBD & \TBD & \TBD & \TBD & \TBD & \TBD \\
    ④ 软奖励        & \TBD & \TBD & \TBD & \TBD & \TBD & \TBD & \TBD & \TBD \\
    ⑤ 连续-无盾     & \TBD & \TBD & \TBD & \TBD & \TBD & \TBD & \TBD & \TBD \\
    ⑥ 连续-有盾     & \TBD & \TBD & \TBD & \TBD & \TBD & \TBD & \TBD & \TBD \\
    ⑦ 消融·都不改   & \TBD & \TBD & \TBD & \TBD & \TBD & \TBD & \TBD & \TBD \\
    ⑧ 消融·只 Beta  & \TBD & \TBD & \TBD & \TBD & \TBD & \TBD & \TBD & \TBD \\
    ⑨ 消融·只对称   & \TBD & \TBD & \TBD & \TBD & \TBD & \TBD & \TBD & \TBD \\
    \bottomrule
  \end{tabular}

  \vspace{4pt}\footnotesize
  \textbf{表注。}官方测试集 600 个场景（与训练集零交集）。训练预算 10{,}158{,}080 步，用 100 个验证场景选最佳存档，全部配置同一规则。到达判据 $=$ 官方 \texttt{is\_reached}（位置 $+$ 朝向 $\pm9.74^\circ$ $+$ 时限）。区间 $=$ 按种子重采样的自助法 95\% 区间；碰撞用 Fisher 精确检验。"练成"判据见 “\nameref{sec:stats}”。
  ⑤⑥ 使用极简配方（无界高斯 / 原文让路入口 / 零塑形）——该配方是无盾连续动作配置唯一合法的形态，故由 ⑤⑥ 得到的“盾的价值”是在连续动作最差形态下测得的。
  转艏增量不可跨配置读作“我们更平顺”：连续动作配置训练时开着直接惩罚该量的罚项、离散动作配置结构上没有；正确读法见图~\ref{fig:ablation}。
  碰撞如实报，与对标的差异按 Fisher 检验陈述，不写“零碰撞”。
\end{table}

\textit{[同表的“仅计达标种子”版本与逐种子原始值：发散种子会污染逐局平均、且污染方向每种配置不同，故两版并报；受篇幅所限，完整两版与逐种子明细放入外部技术报告，正文在此给出关键差值 \TBD。]}

\subsection{消融：两项改进各自贡献多少}
\label{sec:res-ablation}

%  ▶ 消融的定量数据（转艏Δ / 让路违规 / 到达率 / 练成种子数 / 配对 p）随图 3 的原始产物一并给出；
%    受 TRB 图表额度所限（每张折 250 词），消融只出图不出表，完整数值表进外部技术报告。
%    来源：make_头条表.py --ablation ；配对检验来自 stat_检验与区间.py
四条消融配置构成完整的 $2\times2$（有界分布 $\times$ 让路入口），其余配方逐字相同、同种子配对。
定量结果见图~\ref{fig:ablation}：换有界分布把转艏增量降到 \TBD（配对符号检验 $p=$\TBD），
换让路入口把每局让路违规降到 \TBD（$p=$\TBD），两项改进同开为 \TBD。
转艏平顺度唯一站得住的读法就在这组对照里：四种配置开着同一个平顺度罚项，
因此“⑧ 比 ⑦ 平顺 $N$ 倍”只能归因于有界动作分布，而不是额外的奖励项；
跨配置与离散基线比较结构上不成立（我们罚了那个量、离散动作配置拿不到）。
⑦$\to$⑧ 这一级的不可分割伴随（初始探索尺度在加速度轴上差约 2.5 倍）见表~\ref{tab:arms} 表注 (3)。
按 “\nameref{sec:stats}”，本组同样出“仅达标种子”一版（进技术报告）。

\subsection{样本效率与训练可靠性}
\begin{figure}[t]
  \centering
  \figplaceholder{图 2（2$\times$2：学习曲线 / 同种子配对差 / 达标种子数 / 值函数健康度）}
  \caption{样本效率与训练可靠性。（a）九种配置在 100 个验证场景上的到达率学习曲线（每 507{,}904 步一个点，共 20 个点），实线为跨种子中位数、阴影为四分位距。（b）同种子配对的到达率差值（本文方法 $-$ 对标离散动作配置），每条折线一颗种子；折线全部位于零线同侧即为一致优势。（c）“达标种子数”随训练步数的变化（判据见 “\nameref{sec:stats}”）。（d）值函数可解释方差随训练的变化，用于区分“还在爬”与“已经崩”。全部曲线取自训练期落盘的原始数值，未做任何平滑。}
  \label{fig:curves}
\end{figure}
%  ▶ 来源：Paper/正式实验/04_图/make_figures.py → Fig1_sample_efficiency.*
%    数据：各 run 的 progress.json / jsonl 的 trend 与 curves 字段
%  ⚠️ 起飞后待办：该脚本 Fig.1 的 (a)/(c) 两格混了两代实验产物、缺臂时会除零，须先修
%  🔴 不要画"训练期违规曲线"——训练期没采违规量（§4.4），只有那 20 个评估点

\subsection{消融图}
\begin{figure}[t]
  \centering
  \figplaceholder{图 3（1$\times$3：转艏增量 / 让路违规 / 到达率，同种子配对连线）}
  \caption{两项改进的贡献拆解（$2\times2$ 消融，同种子配对）。横轴为四条消融配置，纵轴分别为转艏增量、每局让路违规与到达率。每个点为一颗种子，细线为同种子配对连线，粗线为中位数，误差棒为按种子重采样的自助法 95\% 区间。\textbf{转艏增量的改善应归因于有界动作分布，而非奖励塑形}——四条臂开着同一套塑形（见图~\ref{fig:ablation} 表注）。}
  \label{fig:ablation}
\end{figure}
%  ▶ 来源：同上 → Fig2_ablation.*   ⚠️ 该格目前在缺臂时会静默跳过

\subsection{安全—效率权衡}
\begin{figure}[t]
  \centering
  \figplaceholder{图 4（安全—效率权衡图：到达率 $\times$ 违规，内嵌到达率 $\times$ 碰撞率）}
  \caption{安全—效率权衡。横轴为到达率，纵轴为每局 COLREGs 违规数（对数轴），每个标记为一种配置，误差棒为按种子重采样的自助法 95\% 区间；左上角为理想区。内嵌小图纵轴换为碰撞率。形状编码：$\bullet$ 连续、$\blacktriangle$ 离散、$\circ$ 外部几何基线；实心为带盾、空心为无盾。核心读法：离散动作屏蔽为拿到合规性付出的代价，与连续投影所付代价之比为 \TBD。外部基线的到达率须按 “\nameref{sec:res-ext}” 的限定阅读。}
  \label{fig:tradeoff}
\end{figure}
%  ▶ 来源：同上 → Fig3_tradeoff.*   ⚠️ 该格目前会抛 ValueError

\subsection{盾的行为随训练的演化}
\label{sec:res-shield}
\begin{figure}[t]
  \centering
  \figplaceholder{图 5（2$\times$2：盾介入率 / 动作饱和率两轴 / 盾改写量 / 态势占比）}
  \caption{盾的行为随训练演化（连续动作特有的证据）。（a）盾介入率随训练步数下降，说明策略在 SE-RL 下逐步学会自己产生合规动作。（b）动作饱和率随训练的变化，转艏轴与加速度轴分开画，并把有界分布配置、无界高斯配置与离散动作配置叠在同一张图上——离散网格的 $|\omega|$ 与 $|a|$ 上界恰好等于连续动作箱的半宽，因此两种动作空间共用同一条打满判据、可直接叠图；这是“有界连续动作 vs.\ 离散 $7\times7$ 网格”这一主张的直接证据。（c）投影修正量（按箱半宽归一化）的中位数与 p95 随训练收敛。（d）相遇态势占比随训练的变化，其中紧急态占比是命题 2 作用域的关键数字。}
  \label{fig:shield}
\end{figure}
%  ▶ 🔴 本图需在 Paper/正式实验/04_图/ 新写脚本（现有 make_figures.py 没有这一张）
%    数据：curves 的 roll_* 15 项（盾介入 / 打满舵两轴 / 改写量两轴 / 态势分布）
%    ⚠️ 新写脚本会改掉 BASE_PAPER 树指纹 ⟹ 排到训练跑完再做，做完重算并写回 preflight_formal.sh

\subsection{轨迹定性对比}
\begin{figure}[t]
  \centering
  \figplaceholder{图 6（同一场景 $\times$ 四种方法的轨迹）}
  \caption{同一场景下四种方法的轨迹对比（场景 T-\TBD，种子 s\TBD）。本船轨迹按相遇态势着色（无冲突 / 直航 / 让路 / 紧急），他船轨迹为灰色，圆点间隔 10~s（一个决策步）；虚线圆为命题 1 的远场阈值（约 780~m）。本文方法在让路段做出向右的、幅度超过 $20^\circ$ 的明显转向，而无盾基线从左侧穿过（无碰撞但违规）。本图是定性插图，不承担任何定量声明；场景为人工挑选，种子集合（s0/s1/s2）在跑之前已声明。}
  \label{fig:traj}
\end{figure}
%  ▶ 来源：同上 → Fig4_trajectories.*；轨迹原始数据来自全 600 场景 × 9 臂 × 3 种子的 *_traj.json
%    ⚠️ 该格目前会抛 StopIteration

\subsection{外部几何基线}
\label{sec:res-ext}

%  ▶ 来源：代码/m1_dock_wip/run_baselines_official.py（复用同一套评估代码与同一分母 600）
%    🟢 2026-07-29 已在云端真跑过一趟（此前从未在任何服务器跑过）；PD 两档已出数，VO/CBF 扫参中
\begin{table}[t]
  \caption{外部几何基线对照（不训练、无种子；官方测试集 600）}
  \label{tab:ext}
  \centering\small
  \begin{tabular}{@{}lcccccc@{}}
    \toprule
    方法 & \textbf{到达\%} & \textbf{碰撞\%} & \textbf{违规/局} & \textbf{违规/步} & \textbf{jerk} & \textbf{转艏$\Delta$} \\
    \midrule
    速度障碍 (VO)   & \TBD & \TBD & \TBD & \TBD & \TBD & \TBD \\
    控制障碍函数 (CBF) & \TBD & \TBD & \TBD & \TBD & \TBD & \TBD \\
    PD 控制器       & \TBD & \TBD & \TBD & \TBD & \TBD & \TBD \\
    ① 本文方法      & \multicolumn{6}{c}{（引表~\ref{tab:main}）} \\
    \bottomrule
  \end{tabular}

  \vspace{4pt}\footnotesize
  \textbf{表注。}三条几何基线走无盾路径、不推状态机，因此拿不到相遇态势 $\rho$ $\Rightarrow$ “按态势拆的转艏”一列留空，不补。
  到达率必须带限定：几何控制器的标称律不含终端对准环节，它们位置命中率很高却卡在终端朝向门之外 $\Rightarrow$ 那一列反映的是终端约束类型的差异，不是避碰能力。
  违规数同时报每局与每步两版：它们未到达局会跑满 170 步、我们约 80 步，而直航违规是逐步累计的。
  \emph{不得暗示我们比几何控制器平顺}：正确说法是“我们学出来的策略平顺度与手工调参的几何控制器相当，同时到达率高 \TBD、违规低 \TBD”。
  参数按各基线自己最好的配置报（调参池取自官方训练集，与报数池不重叠），并同时报 RL 动作箱与全物理量程两档。
\end{table}

\subsection{紧急兜底的实证刻画}
\textit{[骨架：① 紧急步占比 \TBD；② 兜底触发率 \TBD；③ 碰撞发生时的归口分布 \TBD——若碰撞全部落在紧急归口，则坐实“紧急态是经验兜底、且落在命题 1 适用域之外”；④ 投影修正量分布 \TBD。来源：重评产物的盾行为字段 $+$ 逐局明细。]}

%==============================================================================
\section{讨论}
\textit{[骨架 · 不含任何超出占位的断言]}

\textbf{投影把“合规”从统计性质变成了逐步性质。}\ 软奖励只能影响策略的期望行为，不能排除任何一个具体动作；屏蔽与投影都能，区别在于屏蔽需要一个可枚举的动作集合。本文的结论是：在连续动作空间里，“可枚举”这个前提可以用“可投影”替代——合规方向恰好是一个半平面，投影到半平面与箱的交集上是一个小规模二次规划，实测单步耗时 \TBD。

\textbf{连续动作的红利落在哪里、不落在哪里。}\ 落在合规转向的精细度上（能精确取到 $\omega_{\mathrm{turn}}$ 下确界）；落在油门轴的平顺度上（结构性，离散只有 7 个油门档，不是训练出来的）；不一定落在到达率上——③$\to$⑤ 这一级很可能是负的，我们照实报（表~\ref{tab:arms} 表注 1）。

\textbf{为什么“两项改进”是连续动作空间特有的问题。}\ 离散动作没有“动作被裁到箱外”这回事，也没有“熵无上界”的通道——这两个病是把 COLREGs 盾搬进连续空间之后新出现的。因此本文的消融不是“又调了两个超参”，而是“补上了迁移过程中新开的两个洞”。

\textbf{与控制理论侧安全滤波器的关系。}\ 反应式控制障碍函数滤波器在 10~s 零阶保持下没有递归可行性证明；本文的可清障集给出了一条路径（“\nameref{sec:prop4}”），但\emph{该结论仍为暂定}，闭环接线与官方口径复算尚未完成，故本文只把它作为设计方向陈述，不作为已兑现的保证。

%==============================================================================
\section{局限}
\begin{enumerate}\itemsep2pt
  \item \textbf{可证明档 $=$ 中。}命题 3 的判据已被验证可靠，但覆盖不完整；命题 4 仍\emph{暂定}；经验零碰撞的支撑目前是单他船的。我们对每条主张都挂满假设，\emph{从不}把安全盾简写成“可证明无碰”。
  \item \textbf{合规是“每让路步”，不是“每步”。}让路态势约占相遇的 $13\text{--}16\%$；紧急态在结构上旁路合规约束；实测每局违规数非零（\TBD）。
  \item \textbf{到达率不列为卖点。}\TBD——若无统计上可辨的优势，如实写“与离散基线相当”。
  \item \textbf{平顺度优势是有条件的。}油门轴上的优势是结构性的；转艏轴上的改善只能以消融对照口径归因于有界动作分布；\emph{不得}暗示我们比手工调参的几何控制器更平顺。
  \item 少数种子的训练不稳定（\TBD）。这是训练不稳定，不是场景难度；奖励塑形对该问题是一条高置信度的经验死路（就我们试过的塑形族与系数而言），但不构成对约束型 MDP 或非 PBRS 终端目标的不可能性证明。
  \item \textbf{COLREGs 覆盖边界与基准局限。}若干条款未实现（与对标工作范围一致）；官方场景库中没有追越场景，追越合规性在本基准上无法验证；基准是双船场景，多他船的方向冲突拿不到实测数据。
  \item \textbf{他船模型与速度上界。}命题 1、3、4 都建立在他船恒速外推上，机动他船会使保证失效；他船速度上界是经验验证的（本基准全池最大 $7.10$~m/s），真实 AIS 数据更快时阈值须相应抬高。
  \item \textbf{报数口径的两个已知偏倚。}最佳存档选择偏倚（期望上抬约 7.5 点，故两版并报）与跨机器跨轮次抖动（连续动作配置约 1.6 个百分点，故强制同一机器、同一轮次评估）。实时性数字随机器而变，只能同机同口径比较。
  \item 离散基线是本文的重新实现（对标工作的海事评估代码未公开），故离散动作配置的种子脆弱性归因于我们的复现，不归因于原作者；软奖励项相对原文有一处数值稳定性偏离，已在方法节声明。
\end{enumerate}

%==============================================================================
\section{结论}
我们把投影式安全盾首次落到海事 COLREGs 的连续动作空间，填上了“连续动作 $\cap$ 可证明方向合规 $\cap$ COLREGs”这个此前空白的三重交集，并精确刻画了这一层能证什么、不能证什么：远场单步无碰严格可证、每让路步方向合规构造式成立、紧急迫近不可避有 可靠性证书；由该证书定义的可清障集的控制不变性已证，据此设计的递归可行性升级\emph{仍为暂定}。在官方场景库上，9 种配置构成的八级单变量递进序列把三组机制逐级拆开（\TBD）。我们如实报告全部局限。

后续工作：把后备机动终端约束接进部署中的盾并验证闭环；构造让路密集 / 对抗性场景源，使兜底与终端约束的价值可被真正检验；多他船的经验验证；把远场判据实现为 $O(1)$ 提前退出以进一步压低单步耗时。

%==============================================================================
\section{致谢}
\textit{[须按投稿要求声明 AI/大模型的使用：用了哪个模型、用于何处。]}

\section{作者贡献}
\textit{[待补]}

\section{利益冲突声明}
作者声明本研究、署名及发表不存在潜在利益冲突。

\section{资助}
\textit{[待补]}

%==============================================================================
\section{参考文献}
\textit{[最后一轮补齐。参考论文的量级是 33 条，TRB 数字顺序引用（按正文首次出现排序）。
锚点：对标工作（动力学 / 状态机 / 指标口径）· 安全盾在环强化学习与 zonotope 投影 · COLREGs 软奖励 ·
深度强化学习报告规范（种子敏感性 / 稳健指标）· CommonOcean 场景库 · COLREGs 公约原文 ·
控制障碍函数与预测式安全滤波 · 速度障碍。]}

\end{document}
